\chapter*{Agradecimientos}

A lo largo de mi vida, he conocido maravillosas personas, que me han aportado todo tipo de experiencias e historias que se han hecho parte de mí, no bastán ni todas las páginas de esta tesis o la novela completa del señor de los anillos para expresar mi más sincero agradecimiento a esas personas.

Esta historia de 4 años (3 y medio para ser específicos) está llena de cuentos como "en la noche termino el proyecto", "claro, puedes coger mi fluido eléctrico", "tenemos una semana para terminar todos los proyectos :)" y un largo etc. Y durante este proceso y a lo largo de mi vida, he conocido maravillosas personas, que me han aportado un sin fín de experiencias que han forjado lo que soy ahora, no bastan ni todas las páginas de la obra de J.K Rowling para expresar mi más profundo agradecimiento.

Gracias a mis padres, Juan Antonio Pedraza Maceda y Rosa Adelaida Romero de la Fe, por ser mi apoyo durante mis primeros años, y siempre estar allí, sin preguntar. Gracias a mi padre por ser mi modelo a seguir, por enseñarme que en la vida no hay imposibles, que todo se puede si te lo propones. Gracias por las croquetas más delicionas que pueden existir en la faz de la tierra jeje.


Muchas gracias a todos y cada uno de los profesores a los que he tenido la oportunidad de recibir sus clases, cada uno de ellos me dejó una impresión que recordaré con nostalgia. Gracias profe Wenny, profe Rodolfo, profe Dariel, profe Martha Dunia, profe Lulu, profe Mayi, profe Tony, profe Lourdes, profe Yolanda, profe Pina, profe Lisandra, profe Vanessa, profe Diansy, profe Adolfo, profe Ernesto. Gracias profe Sonia, tener el placer de recibir sus clases y trabajar con usted durante un tiempo fue un gran placer para mí. Gracias profe Anaisa, por ser un referente pedágogico y por todo lo que he podido aprender de usten referente a Diseño de Software. Gracias profe Raisa por acompañandome durante esta tesis, aconsejandome durante cada paso del proceso y poniendome freno jeje. Gracias profe Fermín.

Clasíficaría mi relación con la Cujae, como de amor-odio, cada fin de semestre es desasperante, pero siendo ansio con añoranza volver a recorrer esos pasillos, sentarme en ese parque y reír con los colegas, gracías CUJAE por estos años que me has dado.